% =====================================================================
% SECTION III — PROBLEM FORMULATION
% File: section3_problem.tex
% =====================================================================

% ------------------------------------------------------------------
\section{Problem Formulation}
\label{sec:problem}
% ------------------------------------------------------------------

Let $\mathcal{G} = (\mathcal{V}, \mathcal{E}, \mathbf{W})$ denote
a directed, typed supply-chain graph where $\mathcal{V}$ is the
set of commodity or product nodes ($|\mathcal{V}|{=}127$ for USGS,
$40$ for SupplyGraph), $\mathcal{E} \subseteq \mathcal{V} \times
\mathcal{V}$ is the set of typed directed edges, and $\mathbf{W}$
encodes edge-type labels and confidence weights.

Each node $v_i$ is associated with a historical demand sequence
$\mathbf{X}_i = [\mathbf{x}^{(t-T+1)}_i, \ldots,
\mathbf{x}^{(t)}_i] \in \mathbb{R}^{T \times d_x}$ over a lookback
window of $T$ timesteps, where $d_x$ is the number of input
features. The forecasting objective is to learn a function

\begin{equation}
  f_\theta: \bigl(\mathbf{X}_i,\,\mathcal{G}\bigr)
  \;\longrightarrow\;
  \hat{y}_i^{(t+1:\,t+h)}
  \label{eq:objective}
\end{equation}

\noindent that maps the joint temporal-structural input to a
$h$-step-ahead demand forecast for each node $v_i$, minimising
a scale-normalised weighted absolute percentage error (WAPE)
across the test set.

We further require the model to produce calibrated uncertainty
estimates $(\hat{y}_{q_{10}},\, \hat{y}_{q_{90}})$ and an
interpretable, per-instance scalar weight $\alpha_i \in (0,1)$
that quantifies the relative contribution of temporal vs.\
structural information to the final prediction for node $v_i$.

\smallskip
\noindent\textbf{Disclosure.} The SupplyGraph dataset covers
222 trading days (2022–-2023). Sequence length $T{=}30$ days for
SupplyGraph and $T{=}3$ years for USGS. Forecast horizon
$h \in \{1, 3, 6, 12\}$ days for SupplyGraph and $h{=}1$ year
for USGS. The small annual USGS dataset is a known limitation;
results should be interpreted in the context of a low-data regime
with only five seed repetitions.
